\VerdierPhysicalPageAnchor{0070}%
\VerdierSemanticAnchor{ch1-2-2-3-continuation}%
are representable:

\VerdierSemanticAnchor{ch1-2-2-3-eq3}
\begin{equation}
  \begin{aligned}
    \left(\int_{\varphi}^{\Pi}Y\right)^{\tau}
      &=\prod_{\varphi(\sigma)=\tau}Y^{\sigma},\\
    d_{\int_{\varphi}^{\Pi}Y}^{\tau,i}
      &\quad\text{defined by the matrix }M(\tau,i,Y),\\[2pt]
    \left(\int_{\varphi}^{\Sigma}Y\right)^{\tau}
      &=\bigoplus_{\varphi(\sigma)=\tau}Y^{\sigma},\\
    d_{\int_{\varphi}^{\Sigma}Y}^{\tau,i}
      &\quad\text{defined by the matrix }M(\tau,i,Y).
  \end{aligned}
  \tag{2.2.3.3}
\end{equation}

\VerdierSemanticAnchor{ch1-2-2-3-pi-extension}
The functor $\int_{\varphi}^{\Pi}$ is called the
$\Pi$-\emph{extension} of the functor $\int_\varphi$.

\VerdierSemanticAnchor{ch1-2-2-3-sigma-extension}
The functor $\int_{\varphi}^{\Sigma}$ is called the
$\Sigma$-\emph{extension} of the functor $\int_\varphi$. These two
extensions will be used below.

\VerdierSemanticAnchor{ch1-2-2-3-natural-map}
There is also a natural morphism of functors
$\int_{\varphi}^{\Sigma}\longrightarrow\int_{\varphi}^{\Pi}$
corresponding to the functorial morphism from the direct sum to the product.

\VerdierSemanticAnchor{ch1-2-2-4}
\noindent\textbf{2.2.4.}\quad
Let
$[n]\xrightarrow{\varphi}[m]\xrightarrow{\psi}[r]$ be two maps and suppose,
to fix ideas, that countable products are representable in $C$. We then have
the relations:

\VerdierSemanticAnchor{ch1-2-2-4-eq1}
\begin{equation}
  \int_{\psi}^{\Pi}\int_{\varphi}^{\Pi}
  =\int_{\psi\varphi}^{\Pi},
  \tag{2.2.4.1}
\end{equation}

\VerdierSemanticAnchor{ch1-2-2-4-eq2}
\begin{equation}
  \int_{\operatorname{id}}^{\Pi}=\operatorname{id}.
  \tag{2.2.4.2}
\end{equation}

\VerdierSemanticAnchor{ch1-2-2-4-editorial-caveat}
The editor finds himself here under the painful obligation of pulling the
usual irritating trick out of the hat. The sign $=$ in formula (2.2.4.1) is
valid only when one has a product functor with suitable associativity
properties, which is never the case in practice.

\VerdierSemanticAnchor{ch1-2-2-4-canonical-isomorphism}
In general, the sign ``$=$'' must be replaced by a canonical isomorphism that
satisfies the usual compatibility properties of the type (1.2.1.3).

\VerdierSemanticAnchor{ch1-2-2-4-sigma-analogue}
Relations analogous to (2.2.4.1) hold for the functor
$\int_{\varphi}^{\Sigma}$.

\VerdierSemanticAnchor{ch1-2-2-4-summable-restriction}
When countable products and sums are not necessarily representable, formula
(2.2.4.1) still holds when one restricts to $\psi\varphi$-summable complexes.

\VerdierSemanticAnchor{ch1-2-2-4-simple-complex}
When $\psi$ is the unique morphism $[n]\longrightarrow[1]$, we use the
notation $\int_n$. For every summable $n$-fold complex $Y$, the complex
$\int_n Y$ is called the \emph{simple complex associated with $Y$}.
